% Equal scale a=3,b=2. Ellipse points show area 3 vs 0;
% the third point lies on the conjugate hyperbola and also gives area 3.
\begin{tikzpicture}[x=.62cm,y=.62cm,>=stealth,font=\small]
\foreach \s/\title in {0/{椭圆：$P=(3,0)$},7.6/{椭圆：$P=(3/\sqrt2,\sqrt2)$},15.2/{共轭双曲线：$P=(0,2)$}}{
\begin{scope}[shift={(\s,0)}]
\node at (0,3.4) {\title};
\draw[gray,dashed] (-3.3,-2.2)--(3.3,2.2) (-3.3,2.2)--(3.3,-2.2);
\node[below left] at (0,0) {$O$};
\end{scope}}
\begin{scope}
\draw[teal!70!black,thick] (0,0) ellipse[x radius=3,y radius=2];
\fill[orange!17] (0,0)--(1.5,1)--(3,0)--(1.5,-1)--cycle;
\draw[orange!85!black,thick] (0,0)--(1.5,1)--(3,0)--(1.5,-1)--cycle;
\fill[red!80!black] (3,0) circle(2pt) node[right] {$P$};
\node at (0,-3) {$S=3$};
\end{scope}
\begin{scope}[shift={(7.6,0)}]
\draw[teal!70!black,thick] (0,0) ellipse[x radius=3,y radius=2];
\draw[orange!85!black,very thick] (0,0)--(2.12132,1.41421);
\fill[red!80!black] (2.12132,1.41421) circle(2pt) node[above left=7pt] {$P=Q$};
\node at (0,-3) {$R=O,\ S=0$（退化）};
\end{scope}
\begin{scope}[shift={(15.2,0)}]
\draw[teal!70!black,thick,domain=-3:3,samples=60] plot (\x,{2*sqrt(1+\x*\x/9)});
\draw[teal!70!black,thick,domain=-3:3,samples=60] plot (\x,{-2*sqrt(1+\x*\x/9)});
\fill[orange!17] (0,0)--(1.5,1)--(0,2)--(-1.5,1)--cycle;
\draw[orange!85!black,thick] (0,0)--(1.5,1)--(0,2)--(-1.5,1)--cycle;
\fill[red!80!black] (0,2) circle(2pt) node[above] {$P$};
\node at (0,-3.7) {$S=3,\quad P\notin H$};
\end{scope}
\end{tikzpicture}
